\chapter{Methods}
\label{sec:Methods}

\section{Dataset}
\label{sec:dataset}

The analysis is based on a catalog of 557 stars associated with the Taurus molecular cloud, combining the membership list of \citet{Luhman2023} with the SigMA clustering algorithm of \citet{Ratzenboeck2023a_SigMA}, both derived from Gaia~DR3 astrometry. For each star the catalog provides the full astrometric solution: equatorial coordinates ($\alpha$, $\delta$), trigonometric parallax $\varpi$ and proper motions in right ascension and declination ($\mu_{\alpha*}$, $\mu_\delta$) together with radial velocities where available, all with their respective uncertainties. As a quality indicator for the astrometric solution serves the renormalized unit weight error (RUWE). All 557 stars possess proper motion measurements, while radial velocities are available for only 180 of them, representing roughly one third of the sample. Photometric data are provided in the form of the Gaia $G$-band magnitude and the colors $G-G_\mathrm{RP}$ and $G_\mathrm{BP}-G_\mathrm{RP}$. Stellar ages, visual extinctions $A_V$, and masses, obtained through isochrone fitting with the \textsc{Chronos} algorithm \citep{Ratzenboeck2023b_Chronos}, are included for a subset of the sample.

The three-dimensional heliocentric Galactic Cartesian positions $(X, Y, Z)$ and corresponding velocity components $(U, V, W)$ are computed from the astrometric and radial-velocity data, with the $X$-axis pointing toward the Galactic center, $Y$ in the direction of Galactic rotation, and $Z$ toward the Galactic north pole.\fxnote{Unklar: Werden die Distanzen direkt aus $1/\varpi$ berechnet, oder aus der Parallaxen-Posterior gezogen (nur das 68.2\%-Kredibilit\"atsintervall behalten), wie es eine fr\"uhere Textversion behauptet hat? Bitte im Code pr\"ufen und hier kurz erg\"anzen -- relevant bei gr\"o{\ss}eren Parallaxenfehlern.}

Prior to inference, the sample is filtered to remove unreliable measurements. Radial velocities outside the range $[-30, 50]~\mathrm{km\,s^{-1}}$ are excluded, as well as stars with uncertainties on their RV measurement greater than XXXX. Such cuts were made only for the RV data because it had far greater uncertainties than the proper motions (+ binaries). Due to the limited number of stars with radial-velocity measurements, this range is set deliberately wide in order to retain as many data points as possible, excluding only the most obvious outliers that are clearly inconsistent with membership in Taurus. Stars with $\mathrm{RUWE} > 1.4$ are additionally excluded from the radial-velocity likelihood, as elevated RUWE values are indicative of unresolved binary systems or a poor astrometric solution that may bias the velocity measurement \citep{Lindegren2021}. Here too, a fixed value was chosen above a position-dependent RUWE threshold derived from the Gaia Unlimited binary selection function \citep{gaiaunlimited}, in order to include as many radial velocities as possible. The resulting catalog spans a spatial volume of roughly $120 \times 75 \times 75~\mathrm{pc}$ around Taurus, discretized on a grid with a cell size of $3~\mathrm{pc}$ per dimension, matching the resolution at which the velocity field is later reconstructed.

\section{Information Field Theory}
\label{sec:ift}

Only a fraction of the stars in the catalog have measured radial velocities, while proper motions, and thus two of the three velocity components, are available for essentially the full sample. Rather than analyzing individual stars or predefined subgroups separately, this thesis combines all of this information into a single, continuous three-dimensional velocity field that describes the motion of Taurus everywhere within the studied volume. However, only at the locations of observed stars does the field yield any meaningful results.

This is a common problem in astrophysics: a continuous physical quantity, here a velocity field, must be reconstructed from a finite, noisy set of measurements. Information field theory (IFT) is a Bayesian framework designed specifically for this kind of problem \citep{Ensslin2023_ITIFT,Ensslin2025}. Instead of treating the unknown field as a small set of parameters, such as the coefficients of a simple analytic function, IFT treats it as the object to be inferred directly: the field is allowed, in principle, to take on any value at any point in space, and its shape is constrained purely by the data and by a set of reasonable prior assumptions.

The key prior assumption used here is that the velocity field is spatially correlated: nearby points are expected to have similar velocities, while the field can vary more or less smoothly depending on the physical scale being considered. This kind of prior is commonly implemented as a Gaussian process, whose correlation properties are described by a power spectrum encoding how much power, i.e.\ variation, the field has on large versus small scales. Given this prior and the observed proper motions and radial velocities, Bayes' theorem in principle yields a posterior probability distribution over all possible velocity fields consistent with the data. In practice, this posterior cannot be written down or sampled from exactly, because the number of unknowns, the field value at every grid point, is far too large. IFT provides the mathematical machinery to approximate this posterior numerically, and to do so in a way that also returns uncertainty estimates for the reconstructed field, not just a single best-fit map.

\section{NIFTy - Numerical Information Field Theory}
\label{sec:nifty}

The numerical implementation of IFT used in this thesis is \textsc{NIFTy}~\citep[\textsc{NIFTy.re};][]{Edenhofer2024}, a Python package built on \textsc{JAX} that extends the original \textsc{NIFTy} framework \citep{niftycl} with ready-to-use building blocks for correlated field models and for the variational inference algorithms needed to approximate the posterior described in Section~\ref{sec:ift}.

Concretely, the three-dimensional velocity field is modeled on a regular grid spanning the volume introduced in Section~\ref{sec:dataset}, with a cell size of $3~\mathrm{pc}$. Each velocity component is represented as a \textsc{NIFTy} correlated field: a constant offset, corresponding to the expected bulk motion of Taurus as a whole, plus a spatially correlated deviation from that offset whose typical smoothness, i.e.\ how much small-scale versus large-scale structure is allowed, is itself inferred from the data rather than fixed by hand.

This general modeling strategy, a correlated field with a flexible, data-driven power spectrum, follows and adapts the approach of \citet{Hutschenreuter2026}, who used the same underlying model and code to reconstruct the 3D velocity field of the Scorpius-Centaurus OB association. With a sample of only about 550 stars in Taurus, splitting the inferred motion into multiple separate kinematic components was not expected to be a meaningful or well-constrained distinction; a single correlated field per velocity component is therefore sufficient to capture the bulk motion and its spatial structure across the region.

Before the field itself is inferred, the noise properties of the stellar velocity measurements are estimated directly from the data, since the nominal Gaia uncertainties do not necessarily capture the full observational error budget. Given the noise model and the correlated-field priors, the posterior distribution over the velocity field is then approximated using \textsc{NIFTy}'s variational inference scheme: the algorithm iteratively refines an approximate posterior by alternating between updating the most probable field configuration and drawing an increasing number of posterior samples around it, in order to capture the local shape of the posterior more accurately as the fit converges. The final result of this procedure is a set of posterior samples of the full 3D velocity field, from which both a mean velocity map and a corresponding uncertainty (standard deviation) map are computed at every grid point, providing the reconstructed velocity field of Taurus together with a quantitative estimate of its reliability.
